\documentclass[12pt,a4paper]{article}
\usepackage[brazil]{babel}
\usepackage[utf8]{inputenc}
\usepackage[T1]{fontenc}
\usepackage{lmodern}
\usepackage{microtype}
\usepackage{geometry}
\usepackage{graphicx}
\usepackage{tabularx}
\usepackage{booktabs}
\usepackage{longtable}
\usepackage{array}
\usepackage{enumitem}
\usepackage{setspace}
\usepackage{caption}
\usepackage{xcolor}
\usepackage{hyperref}
\usepackage{float}
\usepackage{fancyhdr}
\usepackage{amsmath}
\usepackage{url}
\usepackage{xurl}

\geometry{top=2.7cm,bottom=2.4cm,left=3cm,right=2cm}
\onehalfspacing
\setlength{\parindent}{1.25cm}
\setlength{\parskip}{0.2em}
\setlength{\headheight}{14pt}
\setlength{\emergencystretch}{2em}
\setlist[itemize]{leftmargin=1.4cm,itemsep=0.25em,topsep=0.3em}
\setlist[enumerate]{leftmargin=1.4cm,itemsep=0.25em,topsep=0.3em}
\hypersetup{colorlinks=true,linkcolor=blue!55!black,urlcolor=blue!55!black,citecolor=blue!55!black,pdftitle={Atividade XIII: Estratégias de Crescimento, Viabilidade e Coerência — PetCuida},pdfauthor={Grupo 4 — Engenharia de Computação — PUC Minas}}
\captionsetup{font=small,labelfont=bf}
\pagestyle{fancy}
\fancyhf{}
\lhead{\small PetCuida}
\rhead{\small Atividade XIII}
\cfoot{\thepage}
\renewcommand{\headrulewidth}{0.4pt}
\newcolumntype{Y}{>{\raggedright\arraybackslash}X}
\newcolumntype{P}[1]{>{\raggedright\arraybackslash}p{#1}}

\begin{document}
\selectlanguage{brazil}

\begin{titlepage}
\thispagestyle{empty}
\begin{center}
\includegraphics[width=0.19\textwidth]{./assets/brasao.png}\\[0.4cm]
{\large\textbf{PONTIFÍCIA UNIVERSIDADE CATÓLICA DE MINAS GERAIS}}\\[0.15cm]
{\large Curso de Engenharia de Computação}\\
{\large Tópicos III -- Empreendedorismo e Inovação}

\vfill
{\LARGE\textbf{Atividade XIII}}\\[0.35cm]
{\Large Estratégias de crescimento, viabilidade e coerência}\\[0.25cm]
{\large Startup: PetCuida}

\vfill
\textbf{Integrantes}\\[0.15cm]
Gabriel da Silva Cassino\\
Giovanna Naves Ribeiro\\
Júlia Rodrigues Vasconcellos Melo\\
Kathleen Rodrigues Leite Barbosa

\vfill
\textbf{Professor-orientador:} João Carlos Oliveira Caetano\\
\textbf{Professor-tutor:} Paulo Mattos

\vfill
Belo Horizonte\\
2026
\end{center}
\end{titlepage}

\begin{titlepage}
\thispagestyle{empty}
\begin{center}
{\large Gabriel da Silva Cassino}\\
{\large Giovanna Naves Ribeiro}\\
{\large Júlia Rodrigues Vasconcellos Melo}\\
{\large Kathleen Rodrigues Leite Barbosa}

\vfill
{\Large\textbf{Atividade XIII}}\\[0.25cm]
{\large Estratégias de crescimento, viabilidade e coerência}\\[0.25cm]
{\large Startup: PetCuida}
\end{center}
\vfill
\begin{flushright}
\begin{minipage}{8.2cm}
\small\singlespacing
Trabalho acadêmico apresentado à disciplina Tópicos III -- Empreendedorismo e Inovação, do Curso de Engenharia de Computação da Pontifícia Universidade Católica de Minas Gerais, como requisito parcial de avaliação da Atividade XIII.\\[0.35cm]
\textbf{Professor-orientador:} Prof. João Carlos Oliveira Caetano\\
\textbf{Professor-tutor:} Prof. Paulo Mattos
\end{minipage}
\end{flushright}
\vfill
\begin{center}Belo Horizonte\\2026\end{center}
\end{titlepage}

\pagenumbering{roman}
\tableofcontents
\clearpage
\pagenumbering{arabic}

\section{Estratégias de crescimento e viabilidade}

Este documento detalha o plano de crescimento e viabilidade econômica para o PetCuida, conforme os requisitos da Atividade XIII [1]. A estratégia foca na redução de incertezas operacionais e na validação local em Belo Horizonte.

\subsection{Diagnóstico do estágio e sustentabilidade}

O PetCuida é um protótipo em validação. A prioridade é confirmar o retorno de tutores e o valor percebido por parceiros locais. O modelo de receita contempla assinaturas, planos para parceiros, intermediação e patrocínios como hipóteses a testar.

\subsection{Estratégia 1: piloto territorial assistido}

A primeira estratégia é validar o fluxo em um bairro de Belo Horizonte com operação manual assistida (concierge). O objetivo é coletar dados reais de uso antes da automação.

\subsection{Estratégia 2: aquisição por confiança}

Crescer através de conteúdo útil de prevenção e indicações de atores locais, medindo ativação (perfil do pet + lembrete) e retenção no trigésimo dia (D30).

\section{Viabilidade econômica}

A viabilidade será avaliada após a coleta de dados de custos variáveis (suporte, pagamento, infraestrutura) no piloto. A receita recorrente mensal seguirá a fórmula:
\[
\text{Receita mensal} = (N_c \times P_c) + (N_a \times T_a) + S,
\]
onde $N_c$ é o número de clínicas, $P_c$ o preço do plano, $N_a$ o volume de atendimentos intermediados, $T_a$ a taxa líquida e $S$ os patrocínios.

\section{Plano de ação e indicadores}

\begin{longtable}{P{1.8cm}P{4.2cm}P{4.0cm}P{2.7cm}}
\caption{Ciclo de validação proposto.}\label{tab:plano}\\
\toprule
\textbf{Período} & \textbf{Ação} & \textbf{Meta proposta} & \textbf{Responsável} \\
\midrule
\endfirsthead
\toprule
\textbf{Período} & \textbf{Ação} & \textbf{Meta proposta} & \textbf{Responsável} \\
\midrule
\endhead
Sem. 1--2 & Entrevistas e mapeamento local. & 15--20 tutores, 5 clínicas. & Giovanna \\
Sem. 3--4 & Revisão de interface e testes autônomos. & 80\% de sucesso em tarefas. & Kathleen + Gabriel \\
Sem. 5--8 & Piloto concierge em BH. & 40--50 tutores. & Júlia + Equipe \\
\bottomrule
\end{longtable}

\section{Riscos e salvaguardas}

Os riscos de execução incluem baixa oferta local, má interpretação da triagem como diagnóstico e segurança de dados. A resposta no piloto foca em limites clínicos claros e conformidade com o guia da ANPD [2].

\clearpage
\section*{Referências}
\addcontentsline{toc}{section}{Referências}
\small
\begin{enumerate}[label={[\arabic*]},leftmargin=1.3cm,itemsep=0.2em]
\item PONTIFÍCIA UNIVERSIDADE CATÓLICA DE MINAS GERAIS. \textit{XIII -- Estratégias de crescimento, viabilidade e coerência}. Orientação da atividade fornecida em PDF ao grupo, 2026.
\item AUTORIDADE NACIONAL DE PROTEÇÃO DE DADOS (ANPD). \textit{Guia orientativo sobre segurança da informação para agentes de tratamento de pequeno porte}. 2026.
\end{enumerate}
\normalsize

\end{document}
